\section{The self-averaging boundary}
\label{sxn:selfavg}

It is a basic principle of statistics that reliable predictions require samples that are representative of the underlying population. The simplest counterexample is estimating the mean of a heavy-tailed power-law distribution. In such a distribution, a small number of extreme observations can dominate the average, so different finite samples may have very different sample means. A more subtle example arises in complex systems near a critical point. There, even when the ensemble mean is well defined, long-range correlations can keep sample-to-sample fluctuations large, so different large samples do not necessarily become representative of one another. In statistical mechanics, these ideas are formalized by the concept of \emph{Self-Averaging}.
\begin{tcolorbox}[
  marginalbox,
  title=\textbf{Self-averaging versus non-self-averaging}
]
Self-averaging asks if relative sample-to-sample fluctuations $\mathcal{V}_M$ disappear
as the system grows.
\begin{itemize}
\item \(\mathcal{V}_M(\mathcal{O})\to0\): a bound can be formed
\item \(\mathcal{V}_M(\mathcal{O})\not\to0\): no shrinking variance bound
\end{itemize}
\takeaway{The self-averaging boundary is identified here where relative fluctuations $\mathcal{V}_M$ cease to vanish.}
\end{tcolorbox}


Self-averaging means that a macroscopic quantity
measured in one large sample becomes close to its average over many samples as
the system size $M$ grows (i.e concentrates).. 
%
%
%In statistical  physics,, this is frequently stated in terms of the 
%data-averaged partition function $Z$ as
%\begin{equation}
%\left\langle\log\mathcal Z_M\right\rangle_{\xi_M}
%\simeq
%\log\left\langle\mathcal Z_M\right\rangle_{\xi_M} .
%\end{equation}
%This is a strong form of self-averaging; more on this below.
%
Let $\mathcal{O}_M$ be a macroscopic observable measured in a system
of size $M$, and let the mean over $M$ samples $\xi_M$ be
\begin{equation}
\mu_M:=\mathbb{E}[\mathcal{O}_M]=\langle\mathcal{O}\rangle_{\xi_M}
\end{equation}
We say that $\mathcal{O}_M$ is self-averaging when it concentrates around $\mu_M$ as
the system size grows. When the second moment is finite, its relative
sample-to-sample variance is
\begin{equation}
\mathcal V_M(\mathcal O)
:=
\frac{\operatorname{Var}(\mathcal O_M)}
{\left|\mathbb E[\mathcal O_M]\right|^2}
=
\mathbb E\!\left[
\left|
\frac{\mathcal O_M}{\mu_M}-1
\right|^2
\right].
\label{eq:relative_variance}
\end{equation}
A standard sufficient test for self-averaging is
$\mathcal{V}_M(\mathcal{O})\to0$ as $M\to\infty$. When
$\mathcal{V}_M(\mathcal{O})$ does not vanish, relative fluctuations persist and
large samples can remain different.
This is \emph{Non-Self-Averaging} in the relative-variance sense.

\begin{figure}[!b]
\centering
\safeincludegraphics[width=\columnwidth]
{figures/self_averaging_small.png}
\caption{\textbf{Self-averaging and the failure of a shrinking bound.}
Here $\mu_M=\mathbb{E}[\mathcal{O}_M]$. In the self-averaging case, the normalized
observable $\mathcal{O}_M/\mu_M$ concentrates around one and the relative variance
$\mathcal{V}_M(\mathcal{O})$ tends to zero, allowing concentration bounds to tighten.
In the illustrated non-self-averaging case, finite relative deviations retain
nonzero probability as the system grows, so no valid upper bound on that
probability can shrink to zero.}
\label{fig:self_average_schematic}
\end{figure}

In NNs, non-self-averaging arises
from two different mechanisms: localized correlation traps, and 
long-range collective non-vanishing fluctuations.
In either case, a model may describe its training sample well while
performing poorly on out-of-sample test data.

We propose that  the strongest learning occurs at the
boundary between self-averaging and non-self-averaging.
And this is not specific to NNs but a fundamental
property of all related learning algorithms.

But first, it is helpful to understand self-averaging
from different theoretical points of view:
Statistical Learning Theory (SLT), Statistical Mechanics (StatMech),
and Random Matrix Theory (RMT).

\subsection{Statistical Learning view}
The connection to Statistical Learning Theory (SLT) here is qualitive; it
is depicted in Figure~\ref{fig:self_average_schematic}.

For a fixed tolerance $\epsilon>0$, consider the probability that one sample differs
appreciably from the ensemble mean,
\begin{equation}
p_M(\epsilon)
:=
\Pr\!\left\{
\left|
\frac{\mathcal{O}_M}{\mu_M}-1
\right|>
\epsilon
\right\}.
\end{equation}
A useful concentration bound has the form
\begin{equation}
p_M(\epsilon)\leq B_M(\epsilon),
\qquad
B_M(\epsilon)\longrightarrow0
\quad\text{as}\quad M\longrightarrow\infty .
\end{equation}
Such a bound says that a sufficiently large sample is likely to be close to
the ensemble mean.

On the self-averaging side, the relative variance
$\mathcal V_M(\mathcal{O})$ decreases, so the distribution narrows and the bound can
tighten. More data or greater system size then produces a more stable and
representative answer.
If
%$\liminf_{M\to\infty} V_M(O) = 0$,
$\mathcal{V}_M(O) \to 0$
this bound shrinks to zero for every
fixed $\epsilon>0$ since
\begin{equation}
p_M(\epsilon)
\leq
\frac{\mathcal V_M(\mathcal O)}{\epsilon^2}.
\label{eq:relative_chebyshev}
\end{equation}
 A sufficiently large sample is then increasingly likely
to lie close to the ensemble mean.

On the non-self-averaging side, write $D_M:=\mathcal O_M/\mu_M-1$
and $V_M:=\mathcal V_M(\mathcal O)$. Suppose $V_M\geq v_0>0$ and
$\mathbb E[D_M^4]\leq K V_M^2$ for sufficiently large $M$, with fixed
$v_0,K$. Splitting at $|D_M|=\epsilon_0:=\sqrt{v_0/2}$ and applying
Cauchy--Schwarz gives
\[
p_M(\epsilon_0)\geq
\frac{(V_M-\epsilon_0^2)^2}{\mathbb E[D_M^4]}
\geq\frac{1}{4K}>0.
\]
The moment condition excludes variance carried solely by increasingly rare
events. Thus no valid bound $B_M(\epsilon_0)$ can vanish: increasing system
size does not restore concentration.


%The HTSR exponent $\alpha\simeq2$ is the empirical candidate for this
%boundary. Values $\alpha<2$ and correlation traps are treated as warning
%signals that the layer has crossed into the non-self-averaging regime.
\clearpage

\subsection{Statistical Mechanics view}
\label{sxn:selfavg_statmech}

In statistical physics, self-averaging is frequently stated in terms of the
data-averaged partition function $\mathcal Z_M$ as
\begin{equation}
\left\langle\log\mathcal Z_M\right\rangle_{\xi_M}
\simeq
\log\left\langle\mathcal Z_M\right\rangle_{\xi_M}.
\label{eq:selfavg_partition}
\end{equation}
Here $\xi_M$ labels one independently prepared sample of size $M$, and
$\langle\cdot\rangle_{\xi_M}$ denotes an average over many such samples.
Equation~\eqref{eq:selfavg_partition} is a strong form of self-averaging:
a single large sample gives nearly the same free-energy contribution as the
sample average.

To expose the fluctuation correction, define
\begin{equation}
\delta_M
:=
\tfrac{
\mathcal Z_M-\left\langle\mathcal Z_M\right\rangle_{\xi_M}
}{
\left\langle\mathcal Z_M\right\rangle_{\xi_M}
},
\qquad
\left\langle\delta_M\right\rangle_{\xi_M}=0,
\label{eq:relative_partition_fluctuation}
\end{equation}
so that
\begin{equation}
\mathcal V_M(\mathcal Z)
:=
\tfrac{
\operatorname{Var}_{\xi_M}(\mathcal Z_M)
}{
\left\langle\mathcal Z_M\right\rangle_{\xi_M}^{2}
}
=
\left\langle\delta_M^2\right\rangle_{\xi_M}.
\label{eq:relative_variance_partition}
\end{equation}
Whenever the expansion is controlled,
\begin{equation}
\begin{aligned}
\left\langle\log\mathcal Z_M\right\rangle_{\xi_M}
&=
\log\left\langle\mathcal Z_M\right\rangle_{\xi_M}
+
\left\langle\log(1+\delta_M)\right\rangle_{\xi_M}
\\
&=
\log\left\langle\mathcal Z_M\right\rangle_{\xi_M}
-\tfrac12\mathcal V_M(\mathcal Z)
+\tfrac13\left\langle\delta_M^3\right\rangle_{\xi_M}
-\cdots.
\end{aligned}
\label{eq:annealed_quenched_expansion}
\end{equation}
Thus the difference between the two sides of
Eq.~\eqref{eq:selfavg_partition} consists of sample-to-sample fluctuation
terms. If the relative variance and higher normalized fluctuations vanish,
the two expressions agree to leading order. If they remain finite, one large
sample need not represent the full collection of samples, which may degrade generalization.

\paragraph{SETOL and the Spectral Energy.}
Section~\ref{sxn:setol} gives the corresponding effective free energy of layer quality (squared)
generating function $\IZQ$, Eq.~\eqref{eq:setol_q_derivative}. We use this
and ,Eq.~\eqref{eq:setol_Q} to write the
ECS-averaged quality (squared)  $\tilde{Q}^{2}=\overline{Q^2}/M_{\Rell}$,
now averaged over the retained ECS $\Rell$:
\begin{equation}
\tilde{Q}^{2}
=
\frac{1}{M_{\Rell}}\sum_{\tilde{\lambda}_{i}\in\Rell}G_{\Xtilde}(\tilde{\lambda}_{i}),
\;\;
G_{\Xtilde}(\tilde{\lambda})
:=
\int_{\ECSstart}^{\tilde{\lambda}}R_{\Xtilde}(z)\,\dd z,
\label{eq:setol_spectral_functional}
\end{equation}
where $R_{\Xtilde}$ is the $R$-transform of the retained spectrum.
Its local free-cumulant expansion is
\begin{equation}
R_{\mathbf X}(z)
=
\kappa^{(f)}_1
+\kappa^{(f)}_2z
+\kappa^{(f)}_3z^2
+\cdots.
\label{eq:rtransform_cumulants}
\end{equation}
Writing
\begin{equation}
m_p
:=
\frac{1}{M_{\Rell}}\Tr{\Xtilde_{\Rell}^{p}}
=
\int\tilde{\lambda}_{\Rell}^p\widehat\rho_{\Rell}(\tilde{\lambda}_{\Rell})\,\dd\tilde{\lambda}_{\Rell},
\label{eq:spectral_moments}
\end{equation}
To simplify the analysis, let $\ECSstart\simeq 0$; evaluating $\hat{Q}^{2}$ gives
\begin{equation}
\tilde{Q}^2
=
\kappa^{(f)}_1m_1
+\tfrac{1}{2}\kappa^{(f)}_2m_2
+\tfrac{1}{3}\kappa^{(f)}_3m_3
+\cdots 
\label{eq:setol_moment_expansion}
\end{equation}
Since $\kappa^{(f)}_1=m_1$, this gives the ECS-average layer quality as
$ \tilde{Q} \simeq  m_1$.
The leading additive spectral quantity is then
\begin{equation}
\widetilde E_{\Rell}
:=m_1
\simeq
\frac{1}{M_{\Rell}}\Tr{\Xtilde_{\Rell}}
=
\int\tilde{\lambda}_{\Rell}\widehat\rho_{\Rell}(\tilde{\lambda}_{\Rell})\,\dd\tilde{\lambda}_{\Rell}.
\label{eq:trace_law_energy}
\end{equation}
Notice that we choose $\Tr{\X}=M$,
but once we are on the gauge slice and have imposed the trace-log condition,
$\Tr{\Xtilde_{\Rell}}\ne M$ (see~\ref{sxn:meanlog_gauge} below).
This is the \emph{Spectral Energy}; it is the relevant macroscopic observable
for our RG power-counting analysis.

\paragraph{Ising-model analogy and the self-averaging test.}
A similar structure appears in the Ising model. Let $\mathcal M$
denote the total magnetization and let a field $h$ enter through
$\mathcal H(h)=\mathcal H_0-h\mathcal M$. Then $\log\mathcal Z_M(h)$ is the
cumulant-generating function for $\mathcal M$:
\begin{equation}
\log\frac{\mathcal Z_M(h)}{\mathcal Z_M(0)}
=
\sum_{n\geq1}
\frac{(\beta h)^n}{n!}\,\kappa_n(\mathcal M).
\label{eq:ising_cumulant_expansion}
\end{equation}
The first field-dependent term is
$\beta h\,\kappa_1(\mathcal M)=\beta h\langle\mathcal M\rangle$, and the
magnetization density is
\begin{equation}
m_M
:=
\frac{\langle\mathcal M\rangle_{\beta}}{M}
=
\frac{1}{\beta M}
\frac{\partial}{\partial h}\log\mathcal Z_M(h).
\label{eq:ising_magnetization_first_cumulant}
\end{equation}
where the brackets $\langle\rangle_{\beta}$ denote the ordinary thermal average within one fixed sample.
Thus the total magnetization is the first
cumulant, and $m_M$ is its leading additive density generated from the Ising
free energy. In the same sense,
$\widetilde E_{\Rell}=m_1$ is the first additive spectral quantity
generated by the SETOL expansion. It is therefore the first quantity to test
for concentration:
\begin{equation}
\mathcal V_{M_{\Rell}}(\widetilde E_{\Rell})
\underset{?}{\longrightarrow} 0.
\label{eq:energy_selfavg_condition}
\end{equation}
In the disordered Ising model and related random critical systems, the
sample-to-sample relative width of the magnetization and susceptibility does
not necessarily narrow at the critical point. Instead it can remain finite as
the system grows, so these observables are non-self-averaging
\cite{WisemanDomany1995,aharony1996,wiseman1998}.
The analogy here is that if $\widetilde E_{\Rell}$ does not self-average, then the leading SETOL
contribution already depends on the particular training sample or run.
\paragraph{Why criticality defeats averaging.}
The sample-to-sample fluctuations above compare independently prepared
systems; they are not the ordinary thermal fluctuations inside one fixed
sample. Away from criticality, the correlation length $\ell_{\rm corr}$ is
finite, so a $d$-dimensional system of linear size $L$ contains roughly
$(L/\ell_{\rm corr})^d$ nearly independent correlated regions, whose
differences average away. At a disordered critical point,
$\ell_{\rm corr}\sim L$, so the effective number of independent regions can
remain $O(1)$. Increasing the system size then need not reduce relative
fluctuations. This is one reason why magnetization and other singular observables can
remain non-self-averaging at criticality.

For the retained spectral energy,
\begin{equation}
\operatorname{Var}_{\xi_M}(\widetilde E_{\Rell})
=
\frac{1}{M_{\Rell}^{2}}
\sum_{i,j\in\Rell}
\operatorname{Cov}_{\xi_M}(\tilde{\lambda}_{\Rell,i},\tilde{\lambda}_{\Rell,j}).
\label{eq:spectral_energy_covariance}
\end{equation}
If the eigenvalue contributions fluctuate nearly independently, the
off-diagonal terms are small and the average can concentrate. If many modes
fluctuate coherently, the covariance terms add rather than cancel, and the
relative variance can remain finite even when the spectrum contains many
modes. $\widetilde E_{\Rell}$ plays a similar diagnostic role as magnetization
for critical, non-self-averaging behavior.

Correlation traps, however, provide a second, different
route to non-self-averaging, even at criticality; we discuss this in the next subsection.


\clearpage
